\chapter{Data and Cohort Construction}

\section{Data Sources}

The imaging data are drawn from MIMIC-CXR-JPG (version~2.1.0). Structured clinical inputs are derived from MIMIC-IV-ED, specifically emergency department (ED) stay records and triage tables used in cohort linkage and feature construction. The broader MIMIC-IV database is present in the repository for auxiliary pipelines (e.g., laboratory data), but the primary experiments in this thesis rely exclusively on ED-linked chest radiographs and triage features.

The full data processing pipeline is implemented as a sequence of deterministic scripts, including cohort construction, ED linkage, temporal splitting, and downstream table generation. Key steps are executed through scripts such as \texttt{build\_primary\_imaging\_cohort.py}, \texttt{link\_cxr\_to\_edstays.py}, \texttt{build\_final\_ed\_cohort.py}, and \texttt{build\_temporal\_patient\_split.py}. Configuration details for dataset access and local paths are defined in repository configuration files and are not hard-coded.

\section{Cohort Definition}

Cohort construction follows a deterministic pipeline starting from a chest radiograph manifest derived from MIMIC-CXR-JPG. Studies are restricted to frontal views (posteroanterior or anteroposterior), and a single image per study is selected, with posteroanterior priority when multiple candidates are available.

An imaging timestamp ($t_0$) is parsed from study date and time metadata and used as the temporal reference for all downstream processing. Missingness in $t_0$ is tracked during cohort construction; in the retained ED-linked cohort, no studies exhibit missing imaging timestamps.

Radiograph studies are linked to MIMIC-IV-ED emergency department stays using a temporal overlap rule:
\[
\texttt{intime} \leq t_0 \leq \texttt{outtime}.
\]
Only studies with exactly one matched ED stay are retained ($n_{\text{ed\_matches}} = 1$), ensuring a unique and unambiguous clinical context for each imaging study.

Additional integrity filtering removes rows with missing image paths where applicable. The resulting ED-linked cohort contains 81{,}385 studies from 47{,}404 patients and forms the parent population for all downstream label construction and model evaluation.

The filtering process used to derive the final benchmark cohort is summarized in Figure~\ref{fig:cohort_flow}. Starting from the full manifest, the pipeline restricts the data to frontal studies, performs ED linkage under temporal constraints, and retains the labeled subset used for model development and evaluation.

\begin{figure}[htbp]
    \centering
    \begin{tikzpicture}[
        font=\small,
        >=Latex,
        node distance=1.15cm,
        box/.style={
            draw=black!70,
            rounded corners=3pt,
            align=center,
            minimum width=8.2cm,
            minimum height=1.15cm,
            inner sep=6pt
        },
        arrow/.style={->, thick}
    ]

    \node[box] (raw) {
        \textbf{Raw imaging manifest}\\
        MIMIC-CXR-JPG manifest\\
        377{,}110 rows
    };

    \node[box, below=of raw] (frontal) {
        \textbf{Primary frontal imaging cohort}\\
        PA/AP studies only; one image per study with PA-priority\\
        217{,}922 rows
    };

    \node[box, below=of frontal] (ed) {
        \textbf{Final ED-linked cohort}\\
        Temporal linkage: $\texttt{intime} \leq t_0 \leq \texttt{outtime}$\\
        Exactly one matched ED stay per study\\
        81{,}385 studies
    };

    \node[box, below=of ed] (label) {
        \textbf{Labeled evaluation subset}\\
        CheXpert T1 pneumonia label; \texttt{u\_ignore} policy\\
        9{,}137 studies
    };

    \node[box, below=of label] (split) {
        \textbf{Patient-level temporal split}\\
        Train 7{,}132 \quad Validate 930 \quad Test 1{,}075
    };

    \draw[arrow] (raw) -- (frontal);
    \draw[arrow] (frontal) -- (ed);
    \draw[arrow] (ed) -- (label);
    \draw[arrow] (label) -- (split);

    \end{tikzpicture}
    \caption{Cohort derivation for the primary benchmark. Starting from the full MIMIC-CXR-JPG manifest, the pipeline restricts the data to frontal studies, links studies to emergency department encounters using a temporal overlap rule, and retains the CheXpert-labeled \texttt{u\_ignore} subset used for model development and evaluation. Reported counts correspond to generated cohort reports in the pipeline artifacts.}
    \label{fig:cohort_flow}
\end{figure}

\section{Prediction Time and Feature Availability}

Prediction is anchored at the imaging time, defined as the chest radiograph acquisition timestamp ($t_0$). All model inputs are restricted to information available at or before this reference time. This constraint is enforced to align the prediction task with a well-defined clinical moment and to prevent the inclusion of variables influenced by future events.

For the clinical and multimodal baselines, features are limited to emergency department triage data. These include vital signs and scores (temperature, heart rate, respiratory rate, oxygen saturation, systolic and diastolic blood pressure, pain, and acuity), categorical variables (gender, race, arrival transport), and missingness indicators derived for numeric fields after plausibility filtering.

Free-text chief complaint is excluded from the baseline models, and disposition is explicitly removed due to its dependence on downstream clinical decisions and high leakage risk. Feature preprocessing, including scaling, encoding, and imputation, is fit exclusively on the training split within each modelling pipeline to prevent information leakage across splits.

\section{Label Definition}

The primary supervision target is radiographic pneumonia derived from the CheXpert label set (target~T1 in the project definition). Labels are merged into the cohort using study-level identifiers.

In the canonical pipeline run, labels are joined using \texttt{subject\_id} and \texttt{study\_id} as merge keys, reflecting a fallback strategy when image-level identifiers are not consistently available. Consistency checks are applied during the merge process to detect and resolve conflicts.

The primary label policy (\texttt{u\_ignore}) retains only definite positive and negative annotations and excludes uncertain cases. This results in a labeled subset of 9{,}137 studies used for model training and evaluation. An alternative policy (\texttt{u\_zero}) is implemented in the repository but is not used in the primary experiments.

\section{Dataset Splitting}

Data are partitioned using a patient-level temporal split. Each patient is assigned to a single split based on the timing of their earliest imaging study, and all subsequent studies from that patient inherit the same assignment. This ensures that no patient appears in more than one split and that evaluation reflects temporal generalization.

The temporal split is first constructed at the patient level on the full ED-linked cohort and is then propagated to labeled, imaging, and clinical tables using consistent study-level keys. This guarantees that all model variants are trained and evaluated on identical data partitions.

For the primary labeled dataset (\texttt{u\_ignore}), the resulting split contains 7{,}132 training, 930 validation, and 1{,}075 test studies. The combination of patient-level partitioning and strict feature timing constraints is intended to minimize both subject leakage and temporal leakage in model evaluation.